%==============================================================================
% MATH 5080
% Class:        15
% Date:         Oct 20th
%==============================================================================

\subsection{Alternating Series}

\begin{theorem}
    \label{thm:alternating-series-1}
    Let $\{ a_n \} $ be a monotone sequence
    such that $\lim_{n \to \infty}a_n = 0$. 
    Then the series 
    \[
        \sum\limits_{n = 1}^{\infty} (-1)^{n} a_n
    \]
    converges.
\end{theorem}

\begin{example}
    $\sum\limits_{n= 1}^{\infty}  (-1)^n \frac{1}{n} = 1 - \frac{1}{2} + \frac{1}{3} - \frac{1}{4}$
    converges.
\end{example}

\begin{theorem}
    \label{thm:alternating-series-2}
    Let series $\sum\limits_{}^{} a_n$ satisfy:
    \begin{enumerate}
        \item $\lim_{n \to \infty}a_n = 0$
        \item $\exists N > 0$ and a convergent regrouping $\sum\limits_{}^{} b_k$
        s.t. $b_k$ is the sum of $\leq N$ terms of 
        $\{ a_n \}, \forall k \in \mathbb{N}$
    \end{enumerate} 
    Then $\sum\limits_{}^{} a_n$ converges.
\end{theorem}

\begin{proof}[Proof of \ref{thm:alternating-series-2}]
    Let $S_n$ be the $n$-th partial sum of $\sum\limits_{}^{} a_n$,
    and let $S_{n_k}$ be its subsequence that represents 
    the $k$-th partial sum of $\sum\limits_{}^{} b_k$. 
    Since $\sum\limits_{}^{} b_k$ converges, $\exists L \in \mathbb{R}$
    s.t. 
    \[
        \lim_{k \to \infty} S_{n_k} = L
    \]
    \textbf{Claim}: $\lim_{n \to \infty}S_n = L$ 
    \medskip

    Fix any $n \in \mathbb{N}$, then $\exists k \in \mathbb{N} $
    s.t. $n_k \leq n \leq n_{k + 1}$, hence
    \[
        \left| S_n - S_{n - k} \right| 
        \leq 
        \sum\limits_{j = 1}^{N}  \left| a_{n_{k} + j} \right| 
    \]
    Fix any $\ep > 0$, seek $N_0$ s.t. 
    $ \left| S_n - L \right|  < \ep, \forall n > N_0$.
    \medskip

    Let $K_1 > 0$ s.t. $\forall k \geq K_1, \left| S_{n_k} - L \right| < \frac{\ep}{2} $.
    Let $N_1 > 0$ s.t. $ \forall n > N_0, \left| a_n \right| < \frac{\ep}{2N}$.
    \medskip

    Take $N_0 = \max \{ n_{k+1}, N_1 + N \} $
    then $\forall n > N_0$.
    the unique $k$ s.t. $n_k \leq n \leq n_{k + 1}$
    must satisfy $k \geq K_1$.
    Hence $ \left| S_{n_k} - L\right|  < \frac{\ep}{2}$.
    Also, $n > N_0 \geq N_1 + N \implies$,
    $n_k > N_1$, hence 
    \[
        \left| a_{n_k + j} \right| < \frac{\ep}{2N}
    \]

    By triangle inequality,
    \begin{align*}
        \left| S_n - L \right|  & \leq 
        \left| S_n - S_{n_k} \right|  + \left| S_{n_k} - L\right|  \\
        & \leq \sum\limits_{j=1}^{N} |a_{n_k} + j| + |S_{n_k} - L| \\
        & \leq N \cdot \frac{\ep}{2N} + \frac{\ep}{2} = \ep
    \end{align*}
    
\end{proof}
\begin{proof}[Proof of \ref{thm:alternating-series-1}]
    WLOG, assume $\{ a_n \} $ decreasing, 
    then since $\lim_{n \to \infty}a_n = 0$,
    we know that $a_n \geq 0, \forall n$.
    We want to find a regrouping $b_k$ and apply 
    \ref{thm:alternating-series-1}.
    \medskip

    The sequence we are given is 
    \[
        \sum\limits_{n=1}^{\infty} (-1)^{n}a_n
        = - a_1 +a_2 - a_3 + a_4 + \hdots
    \]
    For convenience, we consider 
    \[
        \sum\limits_{n = 1}^{\infty} (-1)^{n-1} a_n
        = a_1 - a-2 + a_3 - a_4 + \hdots
    \]
    Since $a_n$ decreasing, we have
    \begin{align*}
        - \sum\limits_{n = 1}^{\infty} (-1)^{n-1} a_n
        &= a_1 - a-2 + a_3 - a_4 + \hdots \\
        &= (a_1 - a-2) + (a_3 - a_4) + \hdots \\
        &= b_1 + b_2 + \hdots
    \end{align*}
    where $k = 2n$, and
    \[
        b_{k} = (-1)^{(2k-1)-1}a_{2k-1} + (-1)^{2k-1}a_{2k}
    \] 
    Then we verify the two conditions of \ref{thm:alternating-series-2}
    \begin{itemize}
        \item $\lim_{n \to \infty}a_n = 0$, given 
        \item there exists a convergent regrouping $\sum\limits_{k}^{} b_k$
        where each each regrouping is the sum of two terms of $a_n$
        \begin{itemize}
            \item To show $\sum\limits_{}^{} b_k$ convergent, since $b_k$ 
            nonnegative, suffices to show its partial sums are bounded.
            Consider $S_k$, the $k$-th partial sum of $b_k$
        \begin{align*}
            S_k &= \sum\limits_{i = 1}^{k}  b_i \\
            &= (a_1 - a_2) + (a_3 - a_4) + \hdots + (a_{2k-1}, a_{2k}) \\
            &= a_1 - (a_2 - a_3) - (a_4 - a_5) \hdots - (a_{2k-2}- a_{2k-1}) 
            - a_{2k} \\
            & \leq a_1
        \end{align*}
        \end{itemize} 
        Hence we apply \ref{thm:alternating-series-2} and say that 
        $\sum\limits_{}^{} (-1)^{n-1}a_n$ converges.
    \end{itemize} 
\end{proof}

\begin{theorem}[Dirichlet Convergence Test]
    \label{thm:dirichlet-convergence-test}
    Let 
    \begin{itemize}
        \item $\sum\limits_{n=1}^{\infty} a_n$ be a series with
        uniformly bounded partial sums
        \item $\{ b_n \} $ be a monotone sequence s.t. 
        $b_n \to 0$
    \end{itemize} 
    Then 
    $\sum\limits_{n=1}^{\infty} a_n b_n$ converges and 
    \[
    \left| 
        \sum\limits_{n=1}^{\infty} 
        a_nb_n
    \right|  
    \leq A \left| b_1 \right| 
    \]
    
\end{theorem}

\subsection{Products and rearrangement of series}

\begin{definition}[Rearrangement of series]
    Let $\sum\limits_{}^{} a_n$ be a series. 
    \begin{enumerate}
        \item $\sum\limits_{}^{} a_n^{\prime}$ is a \textbf{rearrangement}
        of $\sum\limits_{}^{} a_n$ if there is a bijection 
        $f: \mathbb{N} \to \mathbb{N}$ s.t. 
        $a_n^{\prime} = a_{f(n)}, \forall n \in \mathbb{N}$.
        \item $\sum\limits_{}^{} a_{f(n)}$ is called a \textbf{bounded rearrangement}
        of $\sum\limits_{}^{} a_n$ if $\exists M > 0$ s.t. 
        $|f(n) - n| \leq M, \forall n \in \mathbb{N}$.
    \end{enumerate} 
\end{definition}

\begin{theorem}[
    Series converges iff every
    bounded rearrangement converges
    ]
    Let $\sum\limits_{}^{} a_{f(n)}$ be a bounded rearrangement 
    of $\sum\limits_{}^{} a_n$. Then
    \[
        \sum\limits_{}^{} a_n \text{ converges } 
        \iff
        \sum\limits_{}^{} a_{f(n)} \text{ converges } 
    \]
    and if they converge then they converge to the same limit
    \[
        \sum\limits_{n = 1}^{\infty} a_{n} =
        \sum\limits_{n = 1}^{\infty} a_{f(n)}
    \]
\end{theorem}

\begin{proof}
    Suffices to show 
    \begin{itemize}
        \item $\sum\limits_{}^{} a_n \implies \sum\limits_{}^{} a_{f(n)}$
        \item $\sum\limits_{}^{} a_n = \sum\limits_{}^{} a_{f(n)}$
    \end{itemize} 
    Fix convergent $\sum\limits_{}^{} a_n$ and 
    bounded rearrangement $f$, i.e. $|f(n) - n| \leq M$.
    \medskip

    Since $\sum\limits_{} a_n $ converges, $a_n \to 0$, we get 
    \[
        \lim_{n \to \infty} \sum\limits_{-M}^{M}  |a_{n +j }| = 0
    \]
    \textit{The proof is to fix $\ep$, and pick $N$ s.t. $|a_n| < \frac{\ep}{2M + 1}$}.
    \medskip

    Let $S_n, S_n^{\prime}$ denote the $n$-th partial sum 
    of $\sum\limits_{}^{} a_n, \sum\limits_{}^{} a_{f(n)}$,
    then 
    \[
        \left| S_n - S_n^{\prime} \right|  \leq \sum\limits_{j = -M}^{M} |a_{n + j}|
        \to 0
    \]
    Since $\lim_{n \to \infty}S_n = \sum\limits_{n=1}^{\infty} a_n$ exists,
    \begin{align*}
        \lim_{n \to \infty}S_n^{\prime} = 
        \lim_{n \to \infty}S_n
    \end{align*}

\end{proof}

\begin{theorem}
    [
        Absolute convergence implies
        absolute convergence of any rearrangement
    ]
    Suppose $\sum\limits_{}^{} a_n$ converges absolutely,
    then all of its rearrangement $\sum\limits_{}^{} a_{f(n)}$
    converges absolutely, and 
    \[
        \sum\limits_{n=1}^{\infty} a_n 
        =
        \sum\limits_{n=1}^{\infty} a_{f(n)} 
    \]
\end{theorem}

\begin{proof}
    We first show 
    $\sum\limits_{n = 1}^{\infty} |a_n| = $
    $\sum\limits_{n = 1}^{\infty} |a_{f(n)}|$,
    we then show that 
    $\sum\limits_{n = 1}^{\infty} a_n = $
    $\sum\limits_{n = 1}^{\infty} a_{f(n)}$.
    \medskip

    \textbf{Step 1:}
    $\sum\limits_{n = 1}^{\infty} |a_n| = $
    $\sum\limits_{n = 1}^{\infty} |a_{f(n)}|$,

    Let $S_n^{\prime}$ denote the $n$-th partial sum of
    $\sum\limits_{}^{} |a_{f(n)}|$.
    We have 
    \[
        S_n^{\prime} \leq \sum\limits_{n = 1}^{\infty} |a_n|
    \]
    Since $\sum\limits_{n=1}^{\infty} |a_{f(n)}|$ nonnegative,
    boundedness of partial sum implies $\sum\limits_{}^{} |a_{f(n)}|$
    converges. By symmetry,
    \[
    \sum\limits_{n=1}^{\infty} |a_{n} |=
    \sum\limits_{n=1}^{\infty} |a_{f(n)} |
    \]
    \textbf{Step 2:} now WTS 
    $\sum\limits_{n = 1}^{\infty} a_n = $
    $\sum\limits_{n = 1}^{\infty} a_{f(n)}$.
    Note that if the sequence is nonnegative then we are done. 
    We restrict our attention to sequences with a positive and negative part.
    Define 
    \begin{align*}
        a_n^{+} = \begin{cases}
            a_n, \text{ if } a_n > 0 \\
            0 \text{ if } a_n \geq 0
        \end{cases} \\
        a_n^{-} = \begin{cases}
            -a_n, \text{ if } a_n < 0 \\
            0 \text{ if } a_n \geq 0
        \end{cases}
    \end{align*}
    Hence we have $\forall n$, 
    $a_n^{\pm} \geq 0$.
    By comparison test, $\sum\limits_{}^{} a_n^{\pm}$ converges.
    Now note that 
    \[
        \sum\limits_{n = 1}^{\infty} a_n 
        = 
        \sum\limits_{n = 1}^{\infty} a_n^{+}  -
        \sum\limits_{n = 1}^{\infty} a_n^{-} 
    \]
    Hence 
    \[
        \sum\limits_{n = 1}^{\infty} a_n 
        = 
        \sum\limits_{n = 1}^{\infty} a_n^{+}  -
        \sum\limits_{n = 1}^{\infty} a_n^{-} 
        =
        \sum\limits_{n = 1}^{\infty} a_{f(n)}n^{+}  -
        \sum\limits_{n = 1}^{\infty} a_{f(n)}n^{-} 
    = \sum\limits_{n=1}^{\infty} a_{f(n)}
    \]
\end{proof}

\begin{remark}
    For any series $\sum\limits_{n=1}^{\infty} a_n$,
    write the $n$-th partial sum as $S_n$, then 
    \[
        S_n = S_n^{+} - S_n^{-}
    \]
    and if the limits of $S_n, S_n^{\pm}$ exist,
    \[
    \sum\limits_{n=1}^{\infty} a_n = 
    \lim_{n \to \infty}S_n = 
    \lim_{n \to \infty}S_n^{+} - 
    \lim_{n \to \infty}S_n^{-} 
    \]
\end{remark}


\begin{remark}
    \remind{add in remark}
\end{remark}

\begin{theorem}[Riemann Rearrangement Theorem]
\end{theorem}
\begin{proof}
    \remind{add in proof}
\end{proof}

\begin{theorem}[Matrix product of absolutely convergent series]
\end{theorem}
\begin{proof}
    \remind{add in proof}
\end{proof}

\begin{theorem}[Cauchy product of convergent and absolutely convergent series]
\end{theorem}
\begin{proof}
    \remind{add in proof}
\end{proof}